\section{Persistencia de ruta, multisecciones de especie y conjugaci\'on de M\"obius}
\label{sec:persistencia-ruta-mobius-familias}

Esta secci\'on precisa una distinci\'on que, si se formula de manera demasiado
abreviada, invierte la l\'ogica de la construcci\'on. La ruta no es un dato que
una part\'icula escoge una vez conocida su masa. Tampoco es un dibujo auxiliar
sin contenido causal. Es la historia observable y orientada de una extensi\'on
superviviente. La extensi\'on conserva, adem\'as, la profundidad compatible, el
registro, los acarreos y la memoria que no caben en un corte finito. En este
sentido preciso,
\[
 \boxed{\text{una part\'icula HMT es la persistencia compatible de una ruta
 enriquecida}.}
\]
La frase y la definici\'on por extensiones expresan el mismo objeto a dos
niveles: la ruta es su historia visible; la extensi\'on es la torre completa de
esa historia con todas sus prolongaciones compatibles.

\subsection{Persistencia como secci\'on de un sistema inverso}

Sea \((\mathcal S_m,\rho_{m+1,m})_{m\geq0}\) el sistema de estados que
sobreviven hasta profundidad \(m\), y sea
\[
 x=(x_m)_{m\geq0}\in\Omega_{\rm HMT}^{\rm surv}
 =\varprojlim_m\mathcal S_m.
\]
Las sombras de ruta satisfacen el cuadrado de compatibilidad
\[
 \operatorname{tr}_{m+1,m}\!\left(\operatorname{sh}_{m+1}(x_{m+1})\right)
 =\operatorname{sh}_{m}\!\left(\rho_{m+1,m}(x_{m+1})\right)
 =\operatorname{sh}_{m}(x_m).
\]
Por tanto
\[
 \gamma_x=\bigl(\gamma_x^{[m]}\bigr)_{m\geq0},
 \qquad \gamma_x^{[m]}=\operatorname{sh}_m(x_m),
\]
es una historia compatible, no una secuencia ajustada retrospectivamente.

Sobre cada corte \(\gamma_x^{[m]}\) se dispone del fibrado discreto
\(\mathcal B_{\gamma_x^{[m]}}\). Una secci\'on persistente es una familia
\[
 s=(s_m)_{m\geq0},
 \qquad
 s_m\in\Gamma\!\left(\mathcal B_{\gamma_x^{[m]}}\right),
 \qquad
 \operatorname{res}_{m+1,m}(s_{m+1})=s_m.
\]
Esta ecuaci\'on es el contenido matem\'atico de la persistencia: cada lectura
finita se prolonga sin contradecir las lecturas anteriores. La tupla de
part\'icula del \cref{chap:particula} a\~nade ocupaci\'on,
representaci\'on y propagaci\'on, pero no altera esta procedencia.

\subsection{Especie interna como cociente orbital de la multisección}

En el atlas non\'adico, cada celda orientada \(c=(r,c')\in\mathcal C_{81}\)
porta una familia de ruta
\[
 \mathcal R(c)=
 (n_A,s_C,k_{\Delta_4},\nu_{120},\nu_{270},\tau_{12}).
\]
La imagen de \(\mathcal R\) contiene cincuenta y seis valores. Si se fija la
carta y su orden, el rango \(\kappa\in\{0,1,2,3,4\}\) dentro de cada fibra
produce el levantamiento inyectivo
\[
 c\longmapsto \bigl(\mathcal R(c),\kappa(c)\bigr).
\]
El alfabeto de cinco estados es m\'inimo, pues existen dos fibras de
cardinal cinco. Esta selecci\'on es interna, prospectiva y ciega a masas y a
nombres de part\'iculas.

El censo completo que sostiene esta afirmaci\'on es
\[
 81=45_{\rm leptones}+36_{\rm quarks}
   =25_{\ell^-}+20_{\nu}+20_{d\text{-tipo}}+16_{u\text{-tipo}}.
\]
Las cincuenta y seis fibras de ruta poseen el histograma
\[
 41\times1+10\times2+2\times3+1\times4+2\times5=81.
\]
Por tanto, \(56\) cuenta firmas de ruta, mientras \(81\) cuenta celdas
orientadas; la memoria \(\kappa\) distingue las celdas que una misma firma no
separa.

Para clasificar especies internas no se invierte este mapa. Se toma el
cociente
\[
 q:\mathcal C_{81}\longrightarrow\Sigma_{13},
 \qquad
 q(c)=\bigl(\operatorname{familia}(c),G(c)\bigr),
\]
cuya imagen comprende trece tipos. La salida asociada a un tipo \(y\) es la
multisecci\'on finita del atlas
\[
 \boxed{
 \mathfrak S_y^{\rm atlas}=
 \{(\mathcal R(c),\kappa(c)):q(c)=y\}.}
\]
Sus cardinales son
\[
 1,8,16;\qquad 2,4,6,8;\qquad 2,4,6,8;\qquad4,12,
\]
para leptones cargados, neutrinos y los dos tipos de quark,
respectivamente. La especie es el tipo com\'un de esa fibra; una part\'icula
concreta es una secci\'on persistente realizada dentro de ella. Exigir que el
nombre de una especie elija por s\'i solo una celda destruye precisamente los
grados de orientaci\'on y conjugaci\'on que la fibra conserva.

La elevaci\'on a extensiones es un objeto tipado diferente. Si
\(\pi_{81}:\Omega_{\rm HMT}^{\rm surv}\to\mathcal C_{81}\) denota el lector
celular de la ruta enriquecida declarada, entonces
\[
 \widetilde{\mathfrak S}_y
 =\{x\in\Omega_{\rm HMT}^{\rm surv}:q(\pi_{81}x)=y\}
\]
es la multisecci\'on enriquecida correspondiente. El certificado finito
demuestra el censo y la obstrucci\'on en
\(\mathfrak S_y^{\rm atlas}\); su elevaci\'on es relativa al lector
\(\pi_{81}\), no una identificaci\'on autom\'atica entre una celda y una
historia profunda.

\begin{theorem}[Obstrucci\'on al selector puntual y excepci\'on electr\'onica]
Sea \(\rho_{\rm c}(r,c')=(10-r,10-c')\) la inversi\'on central. Las trece fibras de
\(q\) son \(\rho_{\rm c}\)-estables. La fibra de lept\'on cargado de nivel cero posee
el \'unico punto fijo, \((5,5)\). Ninguna de las otras doce fibras admite
un selector puntual \(\rho_{\rm c}\)-equivariante si \(\rho_{\rm c}\) act\'ua trivialmente
sobre el tipo grueso.
\end{theorem}
\begin{proof}
La ecuaci\'on \(\rho_{\rm c}(r,c')=(r,c')\) equivale a \(r=c'=5\). El censo del atlas
sit\'ua esa celda en la fibra \((\ell^-,G0)\), que tiene cardinal uno. Sea
ahora \(y\neq(\ell^-,G0)\) y supongamos que existe una secci\'on puntual
equivariante \(a:\Sigma_{13}\to\mathcal C_{81}\). Como \(\rho_{\rm c} y=y\), la
equivarianza exigir\'ia
\[
 a(y)=a(\rho_{\rm c} y)=\rho_{\rm c} a(y).
\]
Luego \(a(y)\) tendr\'ia que ser un punto fijo de \(\rho_{\rm c}\), pero el \'unico
punto fijo pertenece a la fibra electr\'onica. Esto es imposible. Por tanto,
en las otras doce fibras la salida natural es una \(\rho_{\rm c}\)-\'orbita o una
multisecci\'on; elegir un representante requiere orientaci\'on o
representaci\'on adicional.
\end{proof}

El electr\'on es excepcional en tres sentidos compatibles: \((5,5)\) es el
centro geom\'etrico \'unico, la \'unica celda de nivel cero y el centro respecto
del cual se forma el ret\'iculo crudo de diferencias
\[
 \Lambda_{\rm raw}^{\Delta}
 =\bigl\langle A_{\rm raw}(c)-A_{\rm raw}(5,5):
 c\in\mathcal C_{81}\bigr\rangle_{\mathbb Z}.
\]
La firma cruda de esa celda no es nula:
\[
 A_{\rm raw}(5,5)=(24,0,12,0,-12),
 \qquad \tau(5,5)=\texttt{+++---+++---}.
\]
El origen \(v_e=(0,0,0,0,0)\) del ret\'iculo de acci\'on es una convenci\'on
relativa al electr\'on y no la firma cruda de la celda central.

El invariante algebraico electr\'onico se eval\'ua mediante
\[
 \mathfrak e_0=
 \frac{\sqrt3}{4}
 \left(e^{\varphi/\pi^2}-22\alpha_{\rm HMT}^{3}\right)
 (1+15\Delta_4),
 \qquad
 \Delta_4=
 \frac{\pi-e\log\pi}{270}\left(1+\frac{\pi}{729}\right),
\]
\[
 \mathfrak e_0
 =0.510998922488066072509870877191349\ldots.
\]
La expresi\'on es adimensional; la carta
\(\varsigma_{u_E}:y\mapsto y\,u_E\) introduce despu\'es una unidad de energ\'ia.
El centro \((5,5)\) y su unicidad son resultados finitos exactos. La selecci\'on
de los coeficientes \(22\) y \(15=3\cdot5\) se deduce, con la estructura de
partida expl\'icita, mediante el complemento
de las cinco regiones de \(\pi\) en \(K_e=\mathbb Z/9\mathbb Z\times\mathbb F_3\):
\(-22=-|K_e\setminus\iota(R_\pi)|\) y
\(+15=|R_\pi\times\mathbb F_3|\). La firma cruda central conserva un tipo
distinto y act\'ua como control independiente de la selecci\'on basal.

La celda \((5,5)\) no debe confundirse con la firma energ\'etica
\(555555\): esta \'ultima pertenece al representante transversal de una de
las cinco regiones de \(\pi\). Son objetos distintos que comparten el d\'igito
central.

En el sector de quarks, el car\'acter residual
\[
 \operatorname{col}(r,c')=r-c'\pmod3
\]
refina la multisecci\'on en tres clases de doce celdas y satisface
\(\operatorname{col}(\rho_{\rm c} c)=-\operatorname{col}(c)\). Esto construye el
tr\'itico interno de color. La representaci\'on f\'isica de \(SU(3)\) es un
lector posterior; no hace falta para demostrar el reparto finito.

\subsection{Estratificación de extensión, ruta, registro y firma}

Los mapas \(q\), \(\mathcal R\), la proyecci\'on celular y el lector de
realizaci\'on distinguen cuatro clasificaciones relacionadas y tipadas:
\begin{enumerate}[label=(\roman*)]
 \item las trece fibras internas familia--nivel de \(q\);
 \item las cincuenta y seis familias de ruta de \(\mathcal R\);
 \item las ochenta y una posiciones y sus trescientas veinticuatro
       orientaciones;
 \item las realizaciones f\'isicas individualizadas por representaci\'on y codominio.
\end{enumerate}
Las cuatro capas se censan antes del contraste metrol\'ogico. Para cada clase,
la comparación requiere la tupla
\[
(\text{posici\'on},\text{fibra},\text{rutas},\text{lectores},
\text{registro},\text{firma},\text{car\'acter},\text{torres},
\text{carta dimensional}).
\]
Una reconstrucci\'on sobre un subconjunto sólo demuestra ese subcaso y no
reemplaza ninguno de los componentes de la tupla.
Dados \(A\), \(C^*\) y los canales \(q_\pm\), la torre global pertenece al
semigrupo infinito
\[
 s_n=q_-^n-q_+^n,
 \qquad
 n\in\langle90,120\rangle
 =\{90,120\}\cup\{30r:r\geq6\}.
\]
Las alturas \(90,120,180,210,240,270,360,420\) son representantes finitos del
semigrupo; la recurrencia anterior prolonga la torre completa.
y el car\'acter adimensional de una firma es
\[
 \chi_{\rm mass}(v)=\exp\!\left(
 n_AA_{\rm rad}+\frac{s_C}{6}C^*_{\rm rad}
 +k\Delta_4+\nu_{120}s_{120}+\nu_{270}(135\Delta_4^2)
 \right).
\]
Los momentos son globales: las especies los ponderan mediante su ruta, su
registro y su operador de fibra. Los lectores \(K_D\) y \(K_\Omega\)
determinan el refinamiento bidireccional antes del contraste, mientras la
tabla integral de especies publica por separado la genealogía interna, la
carta dimensional y la comparación externa.

\subsection{Conjugaci\'on, anti-secci\'on y banda de M\"obius}

Sobre una palabra dodecaf\'asica se define la involuci\'on exacta
\[
 C_M(\tau)_m=-\tau_{11-m},
 \qquad C_M^2=I.
\]
Los funcionales lineales con pesos sim\'etricos son impares bajo \(C_M\),
mientras las correlaciones cuadr\'aticas compatibles con la reversi\'on son
pares. De este modo, orientaci\'on y carga lineal viven en el sector impar,
mientras la memoria cuadr\'atica vive en el sector par. Esto no demuestra que
el car\'acter exponencial completo de masa sea par: si
\(\chi(\tau)=e^{\ell(\tau)}\) y \(\ell\) es impar, entonces
\(\chi(C_M\tau)=\chi(\tau)^{-1}\).

La involuci\'on anterior act\'ua exactamente sobre palabras. Para hablar de
anti-secci\'on se fija adem\'as un levantamiento \(\widetilde C_M\) al fibrado de
secciones que preserve admisibilidad, registro y truncamientos:
\[
 \operatorname{res}_{m+1,m}\widetilde C_M
 =\widetilde C_M\operatorname{res}_{m+1,m},
 \qquad
 \operatorname{sh}(\widetilde C_Mx)=C_M\operatorname{sh}(x).
\]
S\'olo bajo ese levantamiento la anti-secci\'on de \(s\) es
\(\widetilde C_Ms\) sobre la ruta invertida.

La palabra ``M\"obius'' tiene aqu\'i una formulaci\'on relativa precisa. Sea
\(\ell_{\rm ret}\) un lazo de retorno de la ruta enriquecida que genera
\[
 \langle\ell_{\rm ret}\rangle\simeq\pi_1(S^1)\simeq\mathbb Z.
\]
Este grupo de retorno no es la fase finita \(C_9\). Supongamos que su
levantamiento al
registro posee un car\'acter de signo
\[
 \varepsilon_x:\langle\ell_{\rm ret}\rangle\longrightarrow\{\pm1\},
 \qquad \varepsilon_x(\ell_{\rm ret})=-1.
\]
El fibrado de intervalos asociado es
\[
 \mathcal M_x=
 ([0,1]\times[-1,1])/(0,u)\sim(1,-u).
\]
Su primera clase de Stiefel--Whitney es no trivial; por tanto
\(\mathcal M_x\) es una banda de M\"obius. Una vuelta intercambia las dos
hojas locales y dos vueltas restauran la orientaci\'on. Esto no obliga a que
el registro completo vuelva al mismo estado: retorno de orientaci\'on y retorno
de memoria siguen siendo tipos diferentes.

El teorema es relativo al levantamiento, al lazo de retorno y al car\'acter
\(\varepsilon_x\), todos ellos declarados. No nace de observar dos signos.
La lectura de hojas requiere, adem\'as, distinguir la graduaci\'on partida \(R\)
de la conjugaci\'on \(C\) e imponer
\[
 CRC^{-1}=-R,
 \qquad CT(C^*)C^{-1}=T(-C^*).
\]
La lectura part\'icula--antipart\'icula exige finalmente que la representaci\'on
f\'isica invierta los portadores impares, preserve los portadores cuadr\'aticos
pares y entrelace la evoluci\'on y el car\'acter completo de la forma declarada
por la representaci\'on. Bajo esas premisas, part\'icula y antipart\'icula son
las dos orientaciones conjugadas de una misma banda de persistencia.

\begin{theorem}[Descomposici\'on par--impar sobre una cubierta doble]
\label{pf84:A10-md-mobius}
Sea \(p:\widetilde X\to X\) una cubierta principal doble con involuci\'on
\(\iota\). Para toda \(f\in C^0(\widetilde X,\mathbb R)\), los proyectores
\[
 f^+=\frac12(f+\iota^*f),\qquad
 f^-=\frac12(f-\iota^*f)
\]
dan la descomposici\'on \(f=f^++f^-\), con
\[
 f^+\circ\iota=f^+,\qquad f^-\circ\iota=-f^-.
\]
La parte par desciende a una funci\'on sobre \(X\); la parte impar es una
secci\'on de la l\'inea de signo asociada. Para la cubierta de orientaci\'on
de la banda de M\"obius, \(w_1\ne0\) obstruye una secci\'on global continua
y no nula de esa l\'inea.
\end{theorem}

\begin{proof}
Las identidades de los dos proyectores prueban la suma y las paridades. Una
funci\'on invariante es constante en las fibras de \(p\), por lo que
desciende; una funci\'on antiinvariante cumple la ley de transici\'on de la
representaci\'on signo. En la banda de M\"obius la l\'inea asociada es no
trivial y su primera clase de Stiefel--Whitney es no nula, de modo que toda
secci\'on continua global tiene alg\'un cero. La conclusi\'on excluye una
secci\'on global continua \emph{sin ceros}; no niega la existencia de
secciones con ceros.
\end{proof}

Este lema cl\'asico tipa el sector par--impar de la cubierta; no identifica
por s\'i solo la reversi\'on de palabra, la holonom\'ia de retorno, la
graduaci\'on partida, la conjugaci\'on de hojas ni la paridad de una firma de
masa. Esas identificaciones requieren entrelazadores propios.

\subsection{TRIT algebraico y lector temporal}

La realizaci\'on cuadr\'atica declarada asocia al signo visible del TRIT un tipo
algebraico. Para
\(\bar\tau\in\{+1,0,-1\}\), sea
\[
 \mathcal A_{\bar\tau}=
 \mathbb R[J]/(J^2+\bar\tau).
\]
Entonces
\[
 \begin{array}{c|c|c}
 \bar\tau&J^2&\text{r\'egimen}\ \\ \hline
 +1&-1&\text{el\'iptico/complejo}\ \\
 0&0&\text{parab\'olico/nilpotente}\ \\
 -1&+1&\text{hiperb\'olico/partido}.
 \end{array}
\]
Las exponenciales correspondientes son
\[
 e^{tJ}=\cos t+J\sin t,
 \qquad e^{tJ}=1+tJ,
 \qquad e^{tJ}=\cosh t+J\sinh t,
\]
y en el sector partido aparecen los proyectores
\(P_\pm=(1\pm J)/2\). Estas identidades se verifican exactamente en la
realizaci\'on declarada. La
curvatura no procede del signo aislado: aparece cuando APP aporta enlaces y
dos-celdas y se fija una conexi\'on. En la polarizaci\'on mixta suma--producto
publicada, sobre las \(81\) plaquetas, la curvatura orientada satisface
\[
 F(S,S)=F(P,P)=0,
 \qquad F(S,P)=+1,
 \qquad F(P,S)=-1\pmod9.
\]
La identidad de Stokes finita correspondiente controla que el borde de cada
rect\'angulo reproduzca la suma de sus plaquetas.

Al ordenar las transiciones se obtiene el lector temporal tr\'itico:
\[
 \begin{aligned}
 +1&\longmapsto\text{retorno, memoria y pasado},\\
  0&\longmapsto\text{umbral torsional y presente},\\
 -1&\longmapsto\text{apertura bidireccional y futuro}.
 \end{aligned}
\]
Su precursor exacto es la diferencia entre retorno de fase y retorno de
estado:
\[
 g(k+9)=g(k),
 \qquad B_{k+9}\neq B_k\quad\text{en general}.
\]
En la realizaci\'on geom\'etrica correspondiente se habla de sector
esf\'erico--volum\'etrico, umbral plano con registro torsional y sector
hiperb\'olico--areal. Esta correspondencia define un lector de Mec\'anica
Dimensional sujeto al mapa entre el TRIT y la geometr\'ia; no redefine el TRIT
algebraico ni identifica por s\'i sola
un signo con una curvatura del espacio-tiempo. La expresi\'on ``cristal de
tiempo'' nombra esta arquitectura en la terminología HMT; una realizaci\'on
din\'amica de cristal de tiempo exige adem\'as un Hamiltoniano o un protocolo
peri\'odico de excitaci\'on,
respuesta subarm\'onica robusta y estabilidad.

\subsection{Soporte exacto de la lectura ida--retorno}

El intercambio bidireccional no descansa s\'olo en una imagen verbal. Las
coordenadas angulares enteras
\[
 W_+=6n_A+s_C,
 \qquad W_-=6n_A-s_C
\]
se obtienen mediante
\[
 \binom{W_+}{W_-}
 =\begin{pmatrix}6&1\\6&-1\end{pmatrix}
 \binom{n_A}{s_C},
 \qquad
 \operatorname{SNF}\!\begin{pmatrix}6&1\\6&-1\end{pmatrix}
 =\operatorname{diag}(1,12).
\]
La ret\'icula tiene, por tanto, \`indice doce. La involuci\'on
\(s_C\mapsto-s_C\) intercambia \(W_+\leftrightarrow W_-\) y, de forma
equivalente, \(q_+\leftrightarrow q_-\). Este es el sustrato algebraico exacto
de ida y retorno. Su interpretaci\'on como propagaci\'on avanzada y retardada
requiere el operador de evoluci\'on de la representaci\'on f\'isica.

La proyecci\'on par empleada en la lectura biangular satisface, para toda
funci\'on para la que las expresiones est\'en definidas,
\[
 \frac{f(\phi)+f(-\phi)}2=f_{\rm par}(\phi),
 \qquad
 \frac{(\pi-\phi)^2+(\pi+\phi)^2}{2}=\pi^2+\phi^2.
\]
Son identidades exactas de simetrizaci\'on. La lectura en t\'erminos de una
teor\'ia f\'isica del absorbedor pertenece al funtor de propagaci\'on, no a la
identidad algebraica aislada.

\subsection{Periodicidades espinoriales inequivalentes de \texorpdfstring{\(2\pi\)}{2 pi} y \texorpdfstring{\(4\pi\)}{4 pi}}

En el levantamiento espinorial cl\'asico,
\[
 U(\theta)=\exp\!\left(-\frac{i\theta}{2}\sigma_z\right),
 \qquad U(2\pi)=-I,
 \qquad U(4\pi)=I.
\]
La igualdad es exacta en la doble cubierta \(SU(2)\to SO(3)\). Su aplicaci\'on
a una secci\'on HMT requiere especificar la representaci\'on espinorial del
transporte. Una vez fijada, explica por qu\'e una vuelta cambia el signo de la
secci\'on y dos vueltas lo restauran. Si se adoptan las relaciones CAR en el
\'algebra de ocupaci\'on, la exclusi\'on se deduce de ellas; ni CAR ni Pauli se
infieren del signo de \(2\pi\) sin esa estructura adicional.

La interpretación HMT ``hoja areal \(2\pi\), hoja volum\'etrica \(4\pi\)'' es
otro objeto. El apartado~\ref{sec:bisagra-hojas-pi-phi2} formaliza su n\'ucleo
topol\'ogico: una vuelta de la base cierra el borde proyectado en \(2\pi\) y
termina en la hoja conjugada, mientras el lazo cuadrado cierra la cubierta
orientada en \(4\pi\). El teorema es relativo al levantamiento, a la holonom\'ia y al
lector geom\'etrico declarados. Su identificaci\'on adicional con una
representaci\'on espinorial sigue requiriendo el funtor f\'isico y no se deduce
de compartir los mismos factores.

\subsection{Mapa entre multisecciones internas y representaciones físicas}

El objeto part\'icula, la proyecci\'on
extensi\'on--ruta, el selector interno de las 81 celdas, el centro electr\'onico,
las trece multisecciones finitas de familia, su elevaci\'on relativa al lector
celular \(\pi_{81}\), el tr\'itico de color, el registro, la firma, el car\'acter y
las torres forman la composición
\[
\text{extensi\'on}\longrightarrow\text{ruta}\longrightarrow
\text{multisecci\'on}\longrightarrow\text{firma}\longrightarrow
\text{car\'acter}\longrightarrow\text{torres}.
\]

La obligación restante es construir un funtor de
realizaci\'on que empareje ciertas etiquetas f\'isicas individualizadas y
compuestas con las multisecciones enriquecidas pertinentes y reproduzca las
firmas refinadas declaradas sin consultar masas ni firmas objetivo. Para los
estados compuestos, el registro se compone antes de proyectarse a cinco
coordenadas. El dominio del funtor conserva el electr\'on central, el atlas de
familias y la ley interna de car\'acter.

\enlargethispage{6\baselineskip}
\paragraph{Tipado probatorio de persistencia, atlas y familias.}
La persistencia de ruta, las anti-secciones, la lectura bidireccional y el
lector temporal tr\'itico fijan la arquitectura de partida. El sistema
inverso, el selector interno y la separaci\'on de tipos realizan esa
arquitectura mediante morfismos declarados. Los censos del atlas y las
identidades de \(C_M\) sobre palabras se cumplen en todo el sistema finito;
su levantamiento a secciones y la banda de M\"obius se demuestran con
estructura de partida expl\'icita. Los coeficientes electr\'onicos
\((-22,+15)\) se deducen con la misma estructura expl\'icita. Si una
elevaci\'on abre previamente una diferencia metrol\'ogica o una firma objetivo,
su valor es una reconstrucci\'on calibrada.
